%%%%%%%%%%%%%%%%%%%%%%%%%%%%%%%%%%%%%%%%%
% Academic Title Page
% LaTeX Template
% Version 2.0 (17/7/17)
%
% This template was downloaded from:
% http://www.LaTeXTemplates.com
%
% Original author:
% WikiBooks (LaTeX - Title Creation) with modifications by:
% Vel (vel@latextemplates.com)
%
% License:
% CC BY-NC-SA 3.0 (http://creativecommons.org/licenses/by-nc-sa/3.0/)
%
%%%%%%%%%%%%%%%%%%%%%%%%%%%%%%%%%%%%%%%%%

%-------------------------------------------------------------------------------
%	PACKAGES AND OTHER DOCUMENT CONFIGURATIONS
%-------------------------------------------------------------------------------

\documentclass[11pt]{article}

\usepackage[utf8]{inputenc} % Required for inputting international characters
\usepackage[T1]{fontenc}    % Output font encoding for international characters

\usepackage{mathpazo}       % Palatino font

\usepackage{geometry}       % Required for margins
\usepackage{graphicx}       % Required for adding images
\usepackage{hyperref}       % Required for adding hyperlinks
\usepackage{algorithm}      % Required for adding psuedocode
\usepackage{algpseudocode}  % Required for adding psuedocode
\usepackage{minted}         % Required for adding code syntax highlighting


\geometry{                  % Geometry setup
    a4paper,
    bottom=3.0cm,
}

\hypersetup{                % Hyperlink setup
    colorlinks=true,
    linkcolor=black,
    filecolor=magenta,
    urlcolor=cyan
}


\begin{document}

%-------------------------------------------------------------------------------
%	TITLE PAGE
%-------------------------------------------------------------------------------
% Suppresses page number on the title page
% Subsequent page starts page count at 1
\begin{titlepage} 

    % Defines a new command for horizontal lines, change thickness here
	\newcommand{\HRule}{\rule{\linewidth}{0.5mm}} 
	\center % Center everything on the page

	%------------------------------------------------
	%	Main Heading
	%------------------------------------------------
    % Main heading
	\textsc{\LARGE California State University, Los Angeles}\\[1.5cm] 

	%------------------------------------------------
	%	Title
	%------------------------------------------------
	\HRule\\[0.4cm]
    % Title
    {\huge\bfseries Module Level Outcome: \emph{Design, Analysis and
    Application of Algorithms}}\\[0.4cm] \HRule\\[1.5cm]

	%------------------------------------------------
	%	Logo
	%------------------------------------------------
	\vfill\vfill
    % Include a department/university logo
	\includegraphics[width=0.5\textwidth]{../../ReportTemplate/robosub_logo.png}\\[1cm] 

	%------------------------------------------------
	%	Major Heading
	%------------------------------------------------
    % Major heading
	\textsc{\Large RoboSub}\\[0.5cm] 
    % Minor heading
	% \textsc{\large }\\[0.5cm] 

	%------------------------------------------------
	%	Author(s)
	%------------------------------------------------
	{\large\textit{Members}}\\
	Thomas \textsc{Benson},
	David \textsc{Camacho},
	Bailey \textsc{Canham},
	Brandon \textsc{Cao},\\
	Roberto \textsc{Hernandez},
	Andrew \textsc{Heusser},
    Hector \textsc{Mora-Silva},\\
	Bart \textsc{Rando},
	Victor \textsc{Solis}

	%------------------------------------------------
	%	Date
	%------------------------------------------------
    % Position the date 3/4 down the remaining page
	\vfill\vfill\vfill 
    % Date
	{\large\today} 

	%---------------------------------------------------------------------------
    % Push the date up 1/4 of the remaining page
	\vfill 

\end{titlepage}

%-------------------------------------------------------------------------------

%-------------------------------------------------------------------------------
% Table of Contents
%-------------------------------------------------------------------------------
\tableofcontents
\newpage
%-------------------------------------------------------------------------------

%-------------------------------------------------------------------------------
% Content
%-------------------------------------------------------------------------------

\section{Question 5: Convert Sorted Array to Binary Search Tree}
Problem: \url{https://leetcode.com/problems/convert-sorted-array-to-binary-search-tree/}

\subsection{Pseudocode}

% Psuedocode for the solution to the problem

\begin{algorithm}[H]
    \begin{algorithmic}
        \Function{sortedArrayToBST}{$nums$}
            \If{$nums$ is empty}
                \State \Return \{null\}
            \EndIf
            \State $mid \gets \lfloor \frac{len(nums)}{2} \rfloor$
            \State $root \gets \{nums[mid]\}$
            \State $root.left \gets \Call{sortedArrayToBST}{$nums[0:mid-1]$}$
            \State $root.right \gets \Call{sortedArrayToBST}{$nums[mid+1:len(nums)]$}$
            \State \Return $root$
        \EndFunction
    \end{algorithmic}
\end{algorithm}


\subsection{Code}
\subsubsection{C}
\subsubsection{Python}

\begin{minted}{python}
class Solution:
    def sortedArrayToBST(self, nums: List[int]) -> Optional[TreeNode]:
        def constructor(left, right):
            if left > right:
                return None

            midpoint = (left + right) // 2
            root = TreeNode(nums[midpoint])
            root.left = constructor(left, midpoint-1)
            root.right = constructor(midpoint+1, right)
            return root
        return constructor(0, len(nums)-1)
\end{minted}

\subsubsection{TypeScript}

\begin{minted}{typescript}
    const sortedArrayToBST = (nums: number[]): TreeNode | null => {
    if(nums.length === 0) {
        return null;
    }
    const mid: number = Math.floor(nums.length / 2);
    const root: TreeNode = new TreeNode(nums[mid]);
    root.left = sortedArrayToBST(nums.slice(0, mid));
    root.right = sortedArrayToBST(nums.slice(mid + 1));
    return root;
};
\end{minted}

\newpage

\section{Question 6: Binary Tree Preorder Traversal}
Problem: \url{https://leetcode.com/problems/binary-tree-preorder-traversal/}

\subsection{Pseudocode}

% Psuedocode for the iterative solution to the problem.

\begin{algorithm}[H]
    \begin{algorithmic}
        \Function{preorderTraversal}{$root$}
            \State $stack \gets \{\}$
            \State $result \gets \{\}$
            \If{$root$ is null}
                \State \Return $result$
            \EndIf
            \State \Call{push}{$stack, root$}
            \State \While{$stack$ is not empty}
                \State \Call{pop}{$stack, root$}
                \State \Call{append}{$result, root.val$}
                \If{$root.right$ is not null}
                    \State \Call{push}{$stack, root.right$}
                \EndIf
                \If{$root.left$ is not null}
                    \State \Call{push}{$stack, root.left$}
                \EndIf
            \EndWhile
            \State \Return $result$
        \EndFunction
    \end{algorithmic}
\end{algorithm}

% Psuedocode for the recursive solution to the problem.

\begin{algorithm}[H]
    \begin{algorithmic}
        \Function{preorderTraversal}{$root$}
            \State $result \gets \{\}$
            \State \Call{preorder}{$root, result$}
            \State \Return $result$
        \EndFunction
        \Function{preorder}{$root, result$}
            \If{$root$ is null}
                \State \Return
            \EndIf
            \State \Call{append}{$result, root.val$}
            \State \Call{preorder}{$root.left, result$}
            \State \Call{preorder}{$root.right, result$}
        \EndFunction
    \end{algorithmic}
\end{algorithm}

\subsection{Code}
\subsubsection{C}

\begin{minted}{c}

/**
 * Definition for a binary tree node.
 * struct TreeNode {
 *     int val;
 *     struct TreeNode *left;
 *     struct TreeNode *right;
 * };
 */
/**
 * Note: The returned array must be malloced, assume caller calls free().
 */


 void preorder(struct TreeNode* root, int* result, int* index) {
    if (root == NULL) return;

    // append
    result[*index] = root->val;
    (*index)++;

    preorder(root->left, result, index);
    preorder(root->right, result, index);
}

int* preorderTraversal(struct TreeNode* root, int* returnSize) {
    int* result = malloc(100 * sizeof(int));
    int index = 0;

    preorder(root, result, &index);
    *returnSize = index;

    return result;
}

\end{minted}


\subsubsection{Java}

\begin{minted}{java}
import java.util.*;

class Solution {
    public List<Integer> preorderTraversal(TreeNode root) {
        List<Integer> ans = new ArrayList<Integer>();

        Stack<TreeNode> toVisit = new Stack<TreeNode>();

        if (root != null) {
            toVisit.push(root);
        }

        while(!toVisit.empty()) {
            TreeNode hold = new TreeNode();
            hold = toVisit.pop();

            ans.add(hold.val);

            if (hold.right != null) {
                toVisit.push(hold.right);
            }
            if (hold.left != null) {
                toVisit.push(hold.left);
            }
        }

        return ans;
    }
}
\end{minted}

\subsubsection{JavaScript}

\begin{minted}{javascript}

const preorderTraversal = (root) => {
    const result = [];
    const stack = [];
    if (root === null) {
        return result;
    }
    stack.push(root);
    while (stack.length > 0) {
        const node = stack.pop();
        result.push(node.val);
        if (node.right !== null) {
            stack.push(node.right);
        }
        if (node.left !== null) {
            stack.push(node.left);
        }
    }
    return result;
};

\end{minted}

%-------------------------------------------------------------------------------
\end{document}
